\IfFileExists{stacks-project.cls}{%
\documentclass{stacks-project}
}{%
\documentclass{amsart}
}

% For dealing with references we use the comment environment
\usepackage{verbatim}
\newenvironment{reference}{\comment}{\endcomment}
%\newenvironment{reference}{}{}
\newenvironment{slogan}{\comment}{\endcomment}
\newenvironment{history}{\comment}{\endcomment}

% For commutative diagrams we use Xy-pic
\usepackage[all]{xy}

% We use 2cell for 2-commutative diagrams.
\xyoption{2cell}
\UseAllTwocells

% We use multicol for the list of chapters between chapters
\usepackage{multicol}

% This is generally recommended for better output
\usepackage{lmodern}
\usepackage[T1]{fontenc}

% For cross-file-references
\usepackage{xr-hyper}

% Package for hypertext links:
\usepackage{hyperref}

% For any local file, say "hello.tex" you want to link to please
% use \externaldocument[hello-]{hello}
\externaldocument[introduction-]{introduction}
\externaldocument[conventions-]{conventions}
\externaldocument[sets-]{sets}
\externaldocument[categories-]{categories}
\externaldocument[topology-]{topology}
\externaldocument[sheaves-]{sheaves}
\externaldocument[sites-]{sites}
\externaldocument[stacks-]{stacks}
\externaldocument[fields-]{fields}
\externaldocument[algebra-]{algebra}
\externaldocument[brauer-]{brauer}
\externaldocument[homology-]{homology}
\externaldocument[derived-]{derived}
\externaldocument[simplicial-]{simplicial}
\externaldocument[more-algebra-]{more-algebra}
\externaldocument[smoothing-]{smoothing}
\externaldocument[modules-]{modules}
\externaldocument[sites-modules-]{sites-modules}
\externaldocument[injectives-]{injectives}
\externaldocument[cohomology-]{cohomology}
\externaldocument[sites-cohomology-]{sites-cohomology}
\externaldocument[dga-]{dga}
\externaldocument[dpa-]{dpa}
\externaldocument[sdga-]{sdga}
\externaldocument[hypercovering-]{hypercovering}
\externaldocument[schemes-]{schemes}
\externaldocument[constructions-]{constructions}
\externaldocument[properties-]{properties}
\externaldocument[morphisms-]{morphisms}
\externaldocument[coherent-]{coherent}
\externaldocument[divisors-]{divisors}
\externaldocument[limits-]{limits}
\externaldocument[varieties-]{varieties}
\externaldocument[topologies-]{topologies}
\externaldocument[descent-]{descent}
\externaldocument[perfect-]{perfect}
\externaldocument[more-morphisms-]{more-morphisms}
\externaldocument[flat-]{flat}
\externaldocument[groupoids-]{groupoids}
\externaldocument[more-groupoids-]{more-groupoids}
\externaldocument[etale-]{etale}
\externaldocument[chow-]{chow}
\externaldocument[intersection-]{intersection}
\externaldocument[pic-]{pic}
\externaldocument[weil-]{weil}
\externaldocument[adequate-]{adequate}
\externaldocument[dualizing-]{dualizing}
\externaldocument[duality-]{duality}
\externaldocument[discriminant-]{discriminant}
\externaldocument[derham-]{derham}
\externaldocument[local-cohomology-]{local-cohomology}
\externaldocument[algebraization-]{algebraization}
\externaldocument[curves-]{curves}
\externaldocument[resolve-]{resolve}
\externaldocument[models-]{models}
\externaldocument[functors-]{functors}
\externaldocument[equiv-]{equiv}
\externaldocument[pione-]{pione}
\externaldocument[etale-cohomology-]{etale-cohomology}
\externaldocument[proetale-]{proetale}
\externaldocument[relative-cycles-]{relative-cycles}
\externaldocument[more-etale-]{more-etale}
\externaldocument[trace-]{trace}
\externaldocument[crystalline-]{crystalline}
\externaldocument[spaces-]{spaces}
\externaldocument[spaces-properties-]{spaces-properties}
\externaldocument[spaces-morphisms-]{spaces-morphisms}
\externaldocument[decent-spaces-]{decent-spaces}
\externaldocument[spaces-cohomology-]{spaces-cohomology}
\externaldocument[spaces-limits-]{spaces-limits}
\externaldocument[spaces-divisors-]{spaces-divisors}
\externaldocument[spaces-over-fields-]{spaces-over-fields}
\externaldocument[spaces-topologies-]{spaces-topologies}
\externaldocument[spaces-descent-]{spaces-descent}
\externaldocument[spaces-perfect-]{spaces-perfect}
\externaldocument[spaces-more-morphisms-]{spaces-more-morphisms}
\externaldocument[spaces-flat-]{spaces-flat}
\externaldocument[spaces-groupoids-]{spaces-groupoids}
\externaldocument[spaces-more-groupoids-]{spaces-more-groupoids}
\externaldocument[bootstrap-]{bootstrap}
\externaldocument[spaces-pushouts-]{spaces-pushouts}
\externaldocument[spaces-chow-]{spaces-chow}
\externaldocument[groupoids-quotients-]{groupoids-quotients}
\externaldocument[spaces-more-cohomology-]{spaces-more-cohomology}
\externaldocument[spaces-simplicial-]{spaces-simplicial}
\externaldocument[spaces-duality-]{spaces-duality}
\externaldocument[formal-spaces-]{formal-spaces}
\externaldocument[restricted-]{restricted}
\externaldocument[spaces-resolve-]{spaces-resolve}
\externaldocument[formal-defos-]{formal-defos}
\externaldocument[defos-]{defos}
\externaldocument[cotangent-]{cotangent}
\externaldocument[examples-defos-]{examples-defos}
\externaldocument[algebraic-]{algebraic}
\externaldocument[examples-stacks-]{examples-stacks}
\externaldocument[stacks-sheaves-]{stacks-sheaves}
\externaldocument[criteria-]{criteria}
\externaldocument[artin-]{artin}
\externaldocument[quot-]{quot}
\externaldocument[stacks-properties-]{stacks-properties}
\externaldocument[stacks-morphisms-]{stacks-morphisms}
\externaldocument[stacks-limits-]{stacks-limits}
\externaldocument[stacks-cohomology-]{stacks-cohomology}
\externaldocument[stacks-perfect-]{stacks-perfect}
\externaldocument[stacks-introduction-]{stacks-introduction}
\externaldocument[stacks-more-morphisms-]{stacks-more-morphisms}
\externaldocument[stacks-geometry-]{stacks-geometry}
\externaldocument[moduli-]{moduli}
\externaldocument[moduli-curves-]{moduli-curves}
\externaldocument[examples-]{examples}
\externaldocument[exercises-]{exercises}
\externaldocument[guide-]{guide}
\externaldocument[desirables-]{desirables}
\externaldocument[coding-]{coding}
\externaldocument[obsolete-]{obsolete}
\externaldocument[fdl-]{fdl}
\externaldocument[index-]{index}

% Theorem environments.
%
\theoremstyle{plain}
\newtheorem{theorem}[subsection]{Theorem}
\newtheorem{proposition}[subsection]{Proposition}
\newtheorem{lemma}[subsection]{Lemma}

\theoremstyle{definition}
\newtheorem{definition}[subsection]{Definition}
\newtheorem{example}[subsection]{Example}
\newtheorem{exercise}[subsection]{Exercise}
\newtheorem{situation}[subsection]{Situation}

\theoremstyle{remark}
\newtheorem{remark}[subsection]{Remark}
\newtheorem{remarks}[subsection]{Remarks}

\numberwithin{equation}{subsection}

% Macros
%
\def\lim{\mathop{\mathrm{lim}}\nolimits}
\def\colim{\mathop{\mathrm{colim}}\nolimits}
\def\Spec{\mathop{\mathrm{Spec}}}
\def\Hom{\mathop{\mathrm{Hom}}\nolimits}
\def\Ext{\mathop{\mathrm{Ext}}\nolimits}
\def\SheafHom{\mathop{\mathcal{H}\!\mathit{om}}\nolimits}
\def\SheafExt{\mathop{\mathcal{E}\!\mathit{xt}}\nolimits}
\def\Sch{\mathit{Sch}}
\def\Mor{\mathop{\mathrm{Mor}}\nolimits}
\def\Ob{\mathop{\mathrm{Ob}}\nolimits}
\def\Sh{\mathop{\mathit{Sh}}\nolimits}
\def\NL{\mathop{N\!L}\nolimits}
\def\CH{\mathop{\mathrm{CH}}\nolimits}
\def\proetale{{pro\text{-}\acute{e}tale}}
\def\etale{{\acute{e}tale}}
\def\QCoh{\mathit{QCoh}}
\def\Ker{\mathop{\mathrm{Ker}}}
\def\Im{\mathop{\mathrm{Im}}}
\def\Coker{\mathop{\mathrm{Coker}}}
\def\Coim{\mathop{\mathrm{Coim}}}

% Boxtimes
%
\DeclareMathSymbol{\boxtimes}{\mathbin}{AMSa}{"02}

%
% Macros for moduli stacks/spaces
%
\def\QCohstack{\mathcal{QC}\!\mathit{oh}}
\def\Cohstack{\mathcal{C}\!\mathit{oh}}
\def\Spacesstack{\mathcal{S}\!\mathit{paces}}
\def\Quotfunctor{\mathrm{Quot}}
\def\Hilbfunctor{\mathrm{Hilb}}
\def\Curvesstack{\mathcal{C}\!\mathit{urves}}
\def\Polarizedstack{\mathcal{P}\!\mathit{olarized}}
\def\Complexesstack{\mathcal{C}\!\mathit{omplexes}}
% \Pic is the operator that assigns to X its picard group, usage \Pic(X)
% \Picardstack_{X/B} denotes the Picard stack of X over B
% \Picardfunctor_{X/B} denotes the Picard functor of X over B
\def\Pic{\mathop{\mathrm{Pic}}\nolimits}
\def\Picardstack{\mathcal{P}\!\mathit{ic}}
\def\Picardfunctor{\mathrm{Pic}}
\def\Deformationcategory{\mathcal{D}\!\mathit{ef}}
\usepackage{cleveref}
\crefname{originallemma}{original lemma}{original lemmas}
\Crefname{originallemma}{Original lemma}{Original lemmas}
\makeatletter
\newlabel{lemma-completion-taut}{{0GX2}{source}{}{sourcepointer-1}{}}
\newlabel{lemma-completion-taut@cref}{{[originallemma][0][]0GX2}{[1][0][]source}{}{}{}}
\makeatother
\crefname{originallemma}{original lemma}{original lemmas}
\Crefname{originallemma}{Original lemma}{Original lemmas}
\makeatletter
\newlabel{lemma-composition-taut}{{0GX3}{source}{}{sourcepointer-2}{}}
\newlabel{lemma-composition-taut@cref}{{[originallemma][0][]0GX3}{[1][0][]source}{}{}{}}
\makeatother
\crefname{originaldefinition}{original definition}{original definitions}
\Crefname{originaldefinition}{Original definition}{Original definitions}
\makeatletter
\newlabel{more-algebra-definition-topological-ring}{{07E8}{source}{}{sourcepointer-3}{}}
\newlabel{more-algebra-definition-topological-ring@cref}{{[originaldefinition][0][]07E8}{[1][0][]source}{}{}{}}
\makeatother
\crefname{originalsection}{original section}{original sections}
\Crefname{originalsection}{Original section}{Original sections}
\makeatletter
\newlabel{more-algebra-section-topological-ring}{{07E7}{source}{}{sourcepointer-4}{}}
\newlabel{more-algebra-section-topological-ring@cref}{{[originalsection][0][]07E7}{[1][0][]source}{}{}{}}
\makeatother
\title{Equivalent characterisations of weakly pre-adic rings, and the completion of a weakly pre-adic ring}
\author{The Stacks Project Authors}
\date{stacks-project:formal-spaces.tex:local-snapshot-20261012; GNU FDL 1.2 or later}
\begin{document}
\maketitle
\noindent\textit{Editorial note.} All delivered mathematical text is the author's own, copied from the pinned Stacks Project chapter file formal-spaces.tex; no statement or proof step was rewritten, reordered or removed.\par\medskip
\noindent\textit{Editorial note.} The list of equivalent conditions is reproduced with the author's own numbering, including the two additional statements that follow the list rather than being part of the equivalence; the card's obligations keep that separation.\par\medskip
\noindent\textit{Editorial note.} The proof contains the author's abbreviation 'ie.' and the sentence 'The implication (4) implies (3) is trivial.'; both lie inside the counted span and are reproduced verbatim.\par\medskip
\noindent\textit{Editorial note.} The counted spans carry 5 navigation stubs; no bibliography citation occurs in the spans.\par\medskip
\noindent\textit{Editorial note.} The preamble span points at the chapter's own preamble.tex shipped with the frozen source, used verbatim as staged.\par\medskip
\noindent\textit{Editorial note.} The cited definition of a weakly pre-adic ring carries a footnote citing [Gabber-Ramero]; since the chapter ships no bibliography, that entry is attached to this card as bibliography.tex, copied character for character from the Stacks Project's own my.bib, so that the citation resolves without any entry being invented.\par\medskip
\noindent\textit{Allowed standard prerequisites (editorial).}\begin{itemize}
\item Weakly pre-adic and weakly adic rings, and their ideals. Same chapter, Definition definition-weakly-adic, invoked at proof lines 1676-1678 and in the statement. Fixes the first of the four conditions and the notion of an ideal as in that definition.
\item An open ideal in a linearly topologized ring is an ideal of definition only up to containing one, and ideals of definition are open. More on Algebra, Definition more-algebra-definition-topological-ring, invoked at proof lines 1681-1683. Gives that an arbitrary ideal of definition contains a power of the chosen one, which is the step from the third to the fourth formulation.
\item Tautness of continuous ring maps is preserved by composing with the completion, and the completion of a linearly topologized ring is complete. Same chapter, Lemmas lemma-completion-taut and lemma-composition-taut, invoked at proof lines 1692-1694, and More on Algebra, Section more-algebra-section-topological-ring, invoked at proof lines 1696-1697. Produce the taut map needed to see that the completion is weakly pre-adic, and its completeness.
\end{itemize}\medskip
\ifdefined\hypertarget\hypertarget{sourcepointer-1}{}\fi
\noindent\textit{Source reference.} Original lemma 0GX2: sources/stacks-formal-spaces/pinned-original-source.tex, source lines 1342--1342. Pointer outside the counted span, invoked once for context; the load-bearing argument of the counted statement is the author's proof reproduced in full in the proof span.\par\medskip
\ifdefined\hypertarget\hypertarget{sourcepointer-2}{}\fi
\noindent\textit{Source reference.} Original lemma 0GX3: sources/stacks-formal-spaces/pinned-original-source.tex, source lines 1357--1357. Pointer outside the counted span, invoked once for context; the load-bearing argument of the counted statement is the author's proof reproduced in full in the proof span.\par\medskip
\ifdefined\hypertarget\hypertarget{sourcepointer-3}{}\fi
\noindent\textit{Source reference.} Original definition 07E8: sources/stacks-chow/tags.txt, source lines 9438--9438. Pointer outside the counted span, invoked once for context; the load-bearing argument of the counted statement is the author's proof reproduced in full in the proof span.\par\medskip
\ifdefined\hypertarget\hypertarget{sourcepointer-4}{}\fi
\noindent\textit{Source reference.} Original section 07E7: sources/stacks-chow/tags.txt, source lines 9437--9437. Pointer outside the counted span, invoked once for context; the load-bearing argument of the counted statement is the author's proof reproduced in full in the proof span.\par\medskip
\begin{definition}
\label{definition-weakly-adic}
\begin{reference}
\cite[Definition 8.3.8]{Gabber-Ramero}
\end{reference}
Let $A$ be a linearly topologized ring.
\begin{enumerate}
\item We say $A$ is {\it weakly pre-adic}\footnote{In \cite{Gabber-Ramero} the
authors say $A$ is {\it $c$-adic}.} if there exists an ideal
$I \subset A$ such that the closure of $I^n$ is open for all $n \geq 0$
and these closures form a fundamental system of open ideals.
\item We say $A$ is {\it weakly adic} if $A$ is weakly pre-adic
and complete\footnote{By our conventions this includes separated.}.
\end{enumerate}
\end{definition}

\begin{lemma}
\label{lemma-weakly-adic}
Let $A$ be a linearly topologized ring. The following are equivalent
\begin{enumerate}
\item $A$ is weakly pre-adic,
\item there exists a taut continuous ring map $A' \to A$
where $A'$ is a pre-adic topological ring, and
\item $A$ is pre-admissible and there exists an ideal of definition $I$
such that the closure of $I^n$ is open for all $n \geq 1$, and
\item $A$ is pre-admissible and for every ideal of definition $I$
the closure of $I^n$ is open for all $n \geq 1$.
\end{enumerate}
The completion of a weakly pre-adic ring is weakly adic.
If $A$ is weakly adic, then $A$ is admissible and has a countable
fundamental system of open ideals.
\end{lemma}

\begin{proof}
Assume (1). Choose an ideal $I$ such that the closure
of $I^n$ is open for all $n$ and such that these closures form
a fundamental system of open ideals. Denote $A' = A$ endowed with
the $I$-adic topology. Then $A' \to A$ is taut by definition and
we see that (2) holds.

\medskip\noindent
Assume (2). Let $I' \subset A'$ be an ideal of definition.
Denote $I$ the closure of $I'A$. Tautness of $A' \to A$ means
that the closures $I_n$ of $(I')^nA$ are open and form a fundamental system of
open ideals. Thus $I = I_1$ is open and the closures of $I^n$
are equal to $I_n$ and hence open and form a fundamental system of
open ideals. Thus certainly $I$ is an ideal of definition
such that the closure of $I^n$ is open for all $n$. Hence (3) holds.

\medskip\noindent
If $I \subset A$ is as in (3), then $I$ is an ideal as in
Definition \ref{definition-weakly-adic} and we see that (1) holds.
Also, if $I' \subset A$ is any other ideal of definition, then
$I'$ is open (see More on Algebra, Definition
\ref{more-algebra-definition-topological-ring})
and hence contains $I^n$ for some $n \geq 1$.
Thus $(I')^m$ contains $I^{nm}$ for all $m \geq 1$ and we
conclude that the closures of $(I')^m$ are open for all $m$.
In this way we see that (3) implies (4).
The implication (4) $\Rightarrow$ (3) is trivial.

\medskip\noindent
Let $A$ be weakly pre-adic. Choose $A' \to A$ as in (2).
By Lemmas \ref{lemma-completion-taut} and
\ref{lemma-composition-taut}
the composition $A' \to A^\wedge$ is taut. Hence $A^\wedge$
is weakly pre-adic by the equivalence of (2) and (1).
Since the completion of a linearly topologized ring $A$ is complete
(More on Algebra, Section \ref{more-algebra-section-topological-ring})
we see that $A^\wedge$ is weakly adic.

\medskip\noindent
Let $A$ be weakly adic. Then $A$ is complete and
and pre-admissible by (1) $\Rightarrow$ (3) and hence $A$
is admissible. Of course by definition $A$ has a countable
fundamental system of open ideals.
\end{proof}

%% Bibliography retained so that the citation carried by the delivered context span of
%% S5-fsp-weakly-adic (lemma-weakly-adic) resolves.
%%
%% The delivered definition of a weakly pre-adic ring cites [Gabber-Ramero] twice, in a
%% footnote for the terminology and in the definition itself.  The entry below is copied
%% character for character from the pinned Stacks Project bibliography file my.bib (the
%% @ARTICLE{Gabber-Ramero} record), which is the file the Stacks Project itself uses to
%% resolve this key.  Nothing was rewritten, reordered or invented, and the citation key
%% was not changed.
%%
%% Only this one key is cited by the card, and only this one entry is retained.
\begin{thebibliography}{99}

\bibitem[Gabber-Ramero]{Gabber-Ramero}
Gabber, Ofer and Ramero, Lorenzo.
\newblock Foundations for almost ring theory--Release 7.5.
\newblock arXiv preprint math/0409584, 2004.

\end{thebibliography}
\end{document}
